\documentclass[11pt,a4paper]{article}
\usepackage{fontspec}
\setmainfont{DejaVu Serif}[Scale=0.92]
\setsansfont{DejaVu Sans}[Scale=0.92]
\setmonofont{DejaVu Sans Mono}[Scale=0.85]

\usepackage{polyglossia}
\setmainlanguage{french}
\usepackage[margin=2.3cm]{geometry}
\usepackage{amsmath,amssymb,amsfonts,bm}
\usepackage{booktabs,array,longtable}
\usepackage[ruled,vlined,linesnumbered,french,onelanguage]{algorithm2e}
\SetKwInput{KwIn}{Entrées}\SetKwInput{KwOut}{Sortie}
\usepackage{tikz}
\usetikzlibrary{arrows.meta,positioning}
\usepackage[most]{tcolorbox}
\usepackage{xcolor}
\usepackage[hidelinks]{hyperref}
\usepackage{enumitem}
\setlist{itemsep=2pt,topsep=4pt}

\definecolor{bleu}{HTML}{1F5FA8}
\newtcolorbox{exemple}[1][]{colback=blue!3,colframe=bleu!60,fonttitle=\bfseries,title={Exemple chiffré #1},breakable}
\newtcolorbox{defbox}[1]{colback=gray!4,colframe=gray!50,fonttitle=\bfseries,title={#1},breakable}
\newcommand{\R}{\mathbb{R}}
\newcommand{\vect}[1]{\bm{#1}}
\newcommand{\mat}[1]{\mathbf{#1}}
\DeclareMathOperator*{\argmax}{arg\,max}
\DeclareMathOperator{\softmax}{softmax}
\DeclareMathOperator{\idf}{idf}
\DeclareMathOperator{\tf}{tf}

\title{\textbf{Fondements scientifiques et calculs}\\[4pt]
\large claim-coverage-checker : instruction d'un sinistre habitation\\
\normalsize avec rappels pour Insurance Claims Assistance}
\date{Version 5 — novembre 2026}

\begin{document}
\maketitle

\begin{abstract}
\noindent Ce document décrit chaque calcul effectué par l'application, de la lecture des conditions
générales au verdict, et la manière dont il est mesuré. Chaque formule est donnée avec tous ses
paramètres, puis illustrée par un petit exemple calculé sur les données réelles du projet (base
de 999 clauses, classifieur entraîné, polices de démonstration). Le fonctionnement du modèle de
langage exécuté par Ollama est détaillé au niveau matriciel. Les algorithmes sont donnés en
pseudo-code, et chaque section renvoie aux références qui permettent d'aller plus loin.
\end{abstract}

\tableofcontents
\newpage

%=====================================================================
\section{Notations}

Les indices sont fixés pour tout le document.

\begin{center}
\begin{tabular}{>{$}l<{$} p{11cm}}
\toprule
\text{Symbole} & Signification \\
\midrule
i = 1,\dots,m & indice des sinistres (ou des questions) $q_i$ ; $m$ = nombre de sinistres \\
j = 1,\dots,n & indice des clauses (ou des documents) $d_j$ ; $n$ = nombre de clauses ($n = 999$ dans la base) \\
k = 1,\dots,K & indice des familles de garanties (classes) ; $K = 6$ \\
t & un terme (mot, groupe de mots ou n-gramme de caractères) ; $V$ = vocabulaire, $|V|$ sa taille \\
\vect{x}_i \in \R^{|V|} & vecteur (colonne) qui représente le sinistre $q_i$ \\
\mat{X} \in \R^{m\times |V|} & matrice dont la ligne $i$ est $\vect{x}_i^\top$ \\
\mat{W} \in \R^{K\times |V|},\ \vect{b}\in\R^K & poids et biais du classifieur \\
x_{i,t},\ w_{k,t} & coefficients scalaires (ligne $i$ ou $k$, colonne $t$) \\
\mathbb{1}[A] & indicatrice : 1 si la condition $A$ est vraie, 0 sinon \\
\bottomrule
\end{tabular}
\end{center}

\noindent Une \emph{matrice} est un tableau de nombres ; le produit $\mat{A}\mat{B}$ d'une matrice
$p\times r$ par une matrice $r\times s$ donne une matrice $p\times s$ dont le coefficient
$(u,v)$ vaut $\sum_{l=1}^{r} a_{u,l}\,b_{l,v}$. Presque tous les calculs de ce document sont des
produits de ce type.

%=====================================================================
\section{Vue d'ensemble de la chaîne}

L'instruction d'un sinistre $q_i$ sur un contrat suit cinq étapes ; chacune produit l'entrée de la
suivante.

\begin{center}
\begin{tikzpicture}[node distance=5mm and 6mm, every node/.style={font=\small},
  box/.style={draw, rounded corners, fill=blue!5, align=center, minimum height=10mm, text width=2.3cm}]
\node[box] (a) {Description\\ du sinistre};
\node[box, right=of a] (b) {Classifieur\\ (familles $k$)};
\node[box, right=of b] (c) {Recherche\\ BM25 des clauses $d_j$};
\node[box, right=of c] (d) {LLM : 3 questions\\ fermées};
\node[box, below=of d] (e) {Verdict et\\ indemnité};
\node[box, below=of b, text width=4.6cm] (f) {Contrôles factuels (SQL) :\\ contrat, version, garantie, délai, montant};
\draw[-{Latex}] (a)--(b); \draw[-{Latex}] (b)--(c); \draw[-{Latex}] (c)--(d); \draw[-{Latex}] (d)--(e);
\draw[-{Latex}] (b)--(f); \draw[-{Latex}] (f)--(e);
\end{tikzpicture}
\end{center}

%=====================================================================
\section{Préparation des textes}

\begin{defbox}{Mots clés}
\textbf{Normalisation} : mise sous une forme unique (minuscules, accents retirés).
\textbf{Token} : unité de texte après découpage. \textbf{Mot vide} : mot très fréquent sans
contenu (« le », « de »). \textbf{Racinisation} (\emph{stemming}) : réduction d'un mot à une racine
commune à ses variantes.
\end{defbox}

\subsection{Normalisation}
La fonction \texttt{normalize} applique la décomposition Unicode NFKD, supprime les signes
diacritiques combinants, met en minuscules et réduit les espaces. Ainsi
« Dégât » $\to$ « degat ».

\subsection{Termes de recherche et synonymes}
Les termes d'une description sont les tokens $[a\text{-}z]^+$ de plus de deux lettres qui ne sont pas
des mots vides, puis, si une clé du dictionnaire métier figure dans le texte, ses synonymes :
\[
T(q_i) = \{\, t \in \mathrm{tokens}(q_i) : t \notin S,\ |t|>2 \,\} \ \cup \bigcup_{c\,\in\,\text{clés}(q_i)} \mathrm{syn}(c),
\]
où $S$ est la liste des mots vides et $\mathrm{syn}(c)$ les expressions associées à la clé $c$.

\subsection{Racinisation}
Une racine grossière retire le premier suffixe de la liste
(\emph{ements, ement, ations, ation, euses, euse, ees, ee, es, e, s}) si au moins quatre lettres
restent. La requête FTS5 met chaque mot en préfixe (\texttt{racine*}) et chaque expression entre
guillemets, le tout relié par \texttt{OR}.

\begin{exemple}[(vélo, données réelles)]
Description : « Mon vélo a été volé dans le local commun de mon immeuble, qui n'était pas fermé à
clé : suis-je couvert ? »

Tokens retenus : \texttt{velo, vole, local, commun, immeuble, ferme, cle}.
Synonymes ajoutés : \texttt{bicyclette, cycle, mobilite durable, engin de deplacement, vol, locaux
a usage commun, parties communes, local collectif, usage commun, serrure, verrou, fermees a cle,
cles}.

Requête : \texttt{velo* OR vole* OR local* OR commun* OR immeubl* OR ferm* OR cle* OR \dots\ OR
"locaux a usage commun" OR "parties communes" OR \dots}

Le mot de l'assuré « local commun » n'apparaît pas dans les contrats ; c'est le synonyme
« locaux à usage commun » qui retrouve l'exclusion décisive de Sogessur.
\end{exemple}

%=====================================================================
\section{Découpage des conditions générales}

\subsection{Détection de deux colonnes}
Une page à deux colonnes lue ligne par ligne mélange les phrases. L'algorithme cherche une
verticale $x$ entre 40\,\% et 60\,\% de la largeur $L$ de la page. Pour chaque ligne $r$ de la
page ($R$ lignes au total), on note $\mathrm{traverse}(r,x)$ si un mot couvre $x$, $\mathrm{gauche}(r,x)$
si tous les mots sont à gauche et $\mathrm{droite}(r,x)$ s'ils sont tous à droite. La verticale $x$ est
une gouttière si
\[
\sum_r \mathbb{1}[\mathrm{traverse}(r,x)] \le 0{,}05\,R,\quad
\sum_r \mathbb{1}[\mathrm{gauche}(r,x)] \ge 0{,}2\,R,\quad
\sum_r \mathbb{1}[\mathrm{droite}(r,x)] \ge 0{,}2\,R,
\]
et l'on retient celle qui est traversée le moins souvent.

\begin{algorithm}[h]
\caption{Gouttière centrale d'une page}
\KwIn{mots de la page (abscisses $x_0, x_1$), largeur $L$}
\KwOut{abscisse de la gouttière ou \textsc{aucune}}
\lIf{moins de 40 mots}{\Return \textsc{aucune}}
regrouper les mots en lignes $r_1,\dots,r_R$ (même ordonnée)\;
meilleur $\leftarrow$ \textsc{aucune}\;
\For{$x = 0{,}4L$ \KwTo $0{,}6L$ par pas de 2}{
  compter les lignes traversées, à gauche, à droite\;
  \If{les trois conditions sont vraies et $x$ est moins traversée que meilleur}{meilleur $\leftarrow x$}
}
\Return meilleur\;
\end{algorithm}

\subsection{Niveaux de titres}
La taille de police du corps est la plus fréquente : $s_0 = \argmax_s \#\{\text{lignes de taille } s\}$.
Une taille $s$ est une taille de titre si $s \ge s_0 + 0{,}8$ pt, sur au moins deux lignes courtes.
Les tailles triées par ordre décroissant forment un nouveau niveau dès qu'elles s'écartent d'au
moins 1,5 pt du niveau précédent. Résultat : 323 clauses (Baloise), 232 (Ethias) et 444 (Sogessur),
soit $n = 999$.

%=====================================================================
\section{Représentation vectorielle TF-IDF}

\begin{defbox}{Mots clés}
\textbf{TF} (\emph{term frequency}) : nombre d'occurrences d'un terme dans un texte.
\textbf{IDF} (\emph{inverse document frequency}) : poids d'autant plus fort que le terme est rare
dans la collection. \textbf{Norme $\ell_2$} : longueur euclidienne d'un vecteur,
$\|\vect{x}\|_2=\sqrt{\sum_t x_t^2}$.
\end{defbox}

Pour une collection de $m$ textes, scikit-learn calcule, pour le texte $i$ et le terme $t$ :
\[
\tf'_{i,t} = 1 + \ln \tf_{i,t}\ \ (\text{si } \tf_{i,t}>0,\ \text{option \texttt{sublinear\_tf}}),\qquad
\idf_t = \ln\frac{1+m}{1+\mathrm{df}_t} + 1,
\]
\[
x_{i,t} = \frac{\tf'_{i,t}\,\idf_t}{\sqrt{\sum_{t'} \big(\tf'_{i,t'}\,\idf_{t'}\big)^2}},
\]
où $\tf_{i,t}$ est le nombre d'occurrences de $t$ dans le texte $i$ et $\mathrm{df}_t$ le nombre de
textes qui contiennent $t$. Le « +1 » dans la fraction évite une division par zéro ; le « +1 » final
garde un poids positif aux termes présents partout.

\begin{exemple}[(trois phrases)]
Collection ($m=3$) : $q_1$ « vol de vélo », $q_2$ « vélo dans la cave », $q_3$ « dégât des eaux dans la cave ».

\smallskip
\begin{tabular}{lcccc}
\toprule
terme $t$ & $\mathrm{df}_t$ & $\idf_t$ & $x_{1,t}$ & $x_{2,t}$ \\
\midrule
vol & 1 & $\ln(4/2)+1 = 1{,}6931$ & 0,6228 & 0 \\
de  & 1 & 1,6931 & 0,6228 & 0 \\
vélo & 2 & $\ln(4/3)+1 = 1{,}2877$ & 0,4736 & 0,5 \\
cave, dans, la & 2 & 1,2877 & 0 & 0,5 chacun \\
\bottomrule
\end{tabular}

\smallskip
Ligne 1 : vecteur brut $(1{,}6931;\ 1{,}6931;\ 1{,}2877)$, de norme
$\sqrt{2\times1{,}6931^2+1{,}2877^2}=\sqrt{7{,}391}=2{,}7187$ ; d'où
$1{,}6931/2{,}7187=0{,}6228$ et $1{,}2877/2{,}7187=0{,}4736$. Ligne 2 : quatre termes de même idf,
donc $1/\sqrt4=0{,}5$ chacun. Ces valeurs sont celles que produit \texttt{TfidfVectorizer}.
\end{exemple}

Le classifieur concatène deux représentations : les mots et groupes de deux mots (1-2 grammes), et
les n-grammes de caractères de 2 à 5 lettres à l'intérieur des mots (\texttt{char\_wb}), qui
rapprochent « brûlé », « brûlure », « brûlée » :
$\vect{x}_i = \big[\vect{x}_i^{\text{mots}};\ \vect{x}_i^{\text{car.}}\big] \in \R^{6\,153}$.
Ce vecteur est \emph{creux} : la description de l'incendie de cuisine n'a que 209 coefficients non
nuls sur 6\,153.

%=====================================================================
\section{Classifieur des familles de garanties}

\subsection{Régression logistique multinomiale}
\begin{defbox}{Mots clés}
\textbf{Logit} $z_k$ : score non borné de la classe $k$. \textbf{Softmax} : transformation de
scores en probabilités positives de somme 1. \textbf{Entropie croisée} : perte qui pénalise une
faible probabilité donnée à la bonne classe.
\end{defbox}
Pour le sinistre $q_i$ :
\[
\vect{z}_i = \mat{W}\vect{x}_i + \vect{b} \in \R^K,\qquad
p_{i,k} = \softmax(\vect{z}_i)_k = \frac{e^{z_{i,k}}}{\sum_{k'=1}^{K} e^{z_{i,k'}}} .
\]
Les paramètres $(\mat{W},\vect{b})$ minimisent, sur les $m$ descriptions d'entraînement de classes
$y_i$,
\[
\mathcal{L}(\mat{W},\vect{b}) = -\sum_{i=1}^{m} \ln p_{i,y_i} \;+\; \frac{1}{2C}\,\|\mat{W}\|_F^2 ,
\qquad \|\mat{W}\|_F^2=\sum_{k,t} w_{k,t}^2,
\]
avec $C=10$ (plus $C$ est grand, plus la régularisation est faible). Le gradient a une forme
matricielle simple :
$\nabla_{\mat{W}}\mathcal{L} = (\mat{P}-\mat{Y})^\top\mat{X} + \mat{W}/C$, où $\mat{P}\in\R^{m\times K}$
contient les $p_{i,k}$ et $\mat{Y}$ les indicatrices $\mathbb{1}[y_i=k]$ ; il est minimisé par
l'algorithme quasi-Newton L-BFGS.

\begin{exemple}[(incendie de cuisine, modèle réel)]
« Une poêle a pris feu dans la cuisine ; les flammes ont abîmé les meubles et le plafond. »

Logits calculés par le modèle : incendie $5{,}668$ ; dégâts des eaux $0{,}275$ ;
RC $-1{,}288$ ; vol $-1{,}287$ ; événements naturels $-1{,}386$ ; bris de glace $-1{,}983$.

Softmax : $e^{5{,}668} = 289{,}5$ ; somme des exponentielles $= 291{,}7$ ;
$p_{\text{incendie}} = 289{,}5/291{,}7 = 0{,}992$.

Explication : la contribution du terme $t$ à la classe $k$ vaut $x_{i,t}\,w_{k,t}$. Les plus fortes
pour « incendie » : \emph{feu} (0,523), \emph{pris feu} (0,445), \emph{feu dans} (0,405),
\emph{les} (0,346), \emph{flammes} (0,238). La présence de « les » montre une limite : un mot vide
peut porter un poids appris sur les données générées.
\end{exemple}

\subsection{Règle de rejet}
Si $\max_k p_{i,k} < 0{,}60$, le modèle hésite : on retient les deux familles les plus probables, et
la recherche des clauses porte sur les deux.

\subsection{Mesures}
Sur le jeu de test indépendant ($m=60$ descriptions écrites à la main) :
exactitude $54/60 = 90\,\%$ ; log-loss
$-\frac1m\sum_i \ln p_{i,y_i} = 0{,}422$ ; score de Brier
$\frac1m\sum_i\sum_k (p_{i,k}-\mathbb{1}[y_i=k])^2 = 0{,}191$ (0 = parfait).
Le F1 macro moyenne, sur les classes, $F1_k = 2P_kR_k/(P_k+R_k)$, où $P_k$ est la précision et
$R_k$ le rappel de la classe $k$.

%=====================================================================
\section{Recherche des clauses : BM25}

\begin{defbox}{Mots clés}
\textbf{BM25} (\emph{Best Matching 25}) : score de pertinence d'un document pour une requête, issu
du modèle probabiliste de Robertson et Spärck Jones. \textbf{Saturation} : au-delà de quelques
occurrences, un terme répété n'ajoute presque plus rien.
\end{defbox}

Pour une requête $q_i$ et une clause $d_j$ à deux colonnes $c$ (titre de section, texte) :
\[
\mathrm{BM25}(q_i,d_j) = \sum_{c}\ \sum_{t \in T(q_i)} \idf_t \cdot
\frac{f_{t,j,c}\,(k_1+1)}{f_{t,j,c} + k_1\Big(1-b+b\,\dfrac{|d_{j,c}|}{\overline{|d_c|}}\Big)},
\qquad
\idf_t = \ln\frac{n-n_t+0{,}5}{n_t+0{,}5},
\]
où $f_{t,j,c}$ est le nombre d'occurrences de $t$ dans la colonne $c$ de la clause $j$, $|d_{j,c}|$
sa longueur en tokens, $\overline{|d_c|}$ la longueur moyenne de la colonne, $n$ le nombre de clauses,
$n_t$ le nombre de clauses qui contiennent $t$, $k_1=1{,}2$ (saturation) et $b=0{,}75$ (correction
de longueur). SQLite FTS5 renvoie $-\mathrm{BM25}$ (plus petit = meilleur).

\begin{exemple}[(terme « commun », base réelle)]
$n = 999$ clauses ; $n_t = 9$ clauses contiennent « commun ».
$\idf = \ln\frac{999-9+0{,}5}{9+0{,}5} = \ln 104{,}26 = 4{,}6469$.

Meilleure clause : définition « Droit commun » (Sogessur, p.\,12) : $f=1$ dans le texte, longueur
$|d|=45$ tokens, longueur moyenne $\overline{|d|} = 76\,202/999 = 76{,}28$.
\[
4{,}6469\times\frac{1\times2{,}2}{1+1{,}2\,(0{,}25+0{,}75\times45/76{,}28)} = 4{,}6469\times\frac{2{,}2}{1{,}8306} = 5{,}584 .
\]
SQLite donne 5,564 : l'écart de 0,4\,\% vient d'un détail interne de FTS5. Cet exemple montre
aussi une limite : « commun » seul retrouve une définition sans rapport ; ce sont l'expression
« locaux à usage commun » et les quotas qui ramènent l'exclusion utile.
\end{exemple}

\subsection{Composition du contexte}
Le score est multiplié par $1{,}5$ pour les clauses des familles prédites. Le contexte transmis au
modèle est ensuite composé par quotas, pour qu'une exclusion décisive ne soit jamais évincée par
des garanties mieux classées : 2 en-têtes de garantie, 5 garanties ou conditions, 5 exclusions,
4 définitions, soit 16 clauses au plus.

\begin{algorithm}[h]
\caption{Sélection des clauses d'un contrat}
\KwIn{description $q_i$, contrat, familles prédites $F$}
\KwOut{liste ordonnée de clauses}
$T \leftarrow$ termes et synonymes de $q_i$ ; candidats $\leftarrow$ 80 meilleures clauses BM25 du contrat\;
\ForEach{candidat $d_j$}{\lIf{une famille de $d_j$ est dans $F$}{score $\leftarrow 1{,}5\times$score}}
trier les candidats par score décroissant\;
ajouter les en-têtes de garantie des familles $F$, puis les définitions des termes de $T$\;
\ForEach{type $\in$ \{garantie/condition, exclusion, définition\}}{
  ajouter les meilleurs candidats de ce type jusqu'au quota}
\Return clauses triées : définitions, garanties et conditions, exclusions\;
\end{algorithm}

%=====================================================================
\section{Le modèle de langage dans Ollama}

\begin{defbox}{Mots clés}
\textbf{LLM} (\emph{large language model}) : réseau de neurones qui estime la probabilité du token
suivant sachant les précédents. \textbf{Transformer} : architecture fondée sur l'attention.
\textbf{Ollama} : serveur local qui télécharge, charge et exécute un modèle et expose une API HTTP
(\texttt{/api/generate}) ; il s'appuie sur la bibliothèque llama.cpp et le format de fichier
GGUF. \textbf{Quantification} : stockage des poids sur peu de bits.
\end{defbox}

\subsection{Le modèle utilisé : \texttt{llama3.2:3b}}
Caractéristiques de Llama 3.2 3B (fichier de configuration publié) :

\begin{center}
\begin{tabular}{lr|lr}
\toprule
dimension cachée $d$ & 3\,072 & couches $L$ & 28 \\
têtes d'attention (requêtes) $H$ & 24 & têtes clé/valeur $H_{kv}$ & 8 \\
dimension d'une tête $d_h$ & 128 & dimension intermédiaire $d_{ff}$ & 8\,192 \\
vocabulaire $|V|$ & 128\,256 & contexte maximal & 131\,072 tokens \\
\bottomrule
\end{tabular}
\end{center}

\begin{exemple}[(nombre de paramètres)]
Matrice d'embedding $\mat{E}\in\R^{|V|\times d}$ : $128\,256\times3\,072 = 394{,}0$ M.

Une couche : attention $\mat{W}_Q,\mat{W}_O$ ($3\,072\times3\,072$ chacune) et $\mat{W}_K,\mat{W}_V$
($3\,072\times1\,024$ chacune, car $H_{kv}d_h = 8\times128$) : $25{,}2$ M ; réseau
\emph{feed-forward} (3 matrices $3\,072\times8\,192$) : $75{,}5$ M ; total 100,7 M.

Total : $394{,}0 + 28\times100{,}7 = 3{,}21$ milliards (la matrice de sortie est partagée avec
$\mat{E}$). Quantifiés à environ 4,85 bits par poids (format Q4\_K\_M), ils occupent
$3{,}21\times10^9\times4{,}85/8 \approx 1{,}95$ Go : c'est la taille du fichier téléchargé.
\end{exemple}

\subsection{Tokenisation BPE}
Le texte est découpé en tokens par un algorithme BPE (\emph{byte pair encoding}) : partant des
octets, on fusionne de façon répétée la paire adjacente la plus fréquente du corpus
d'entraînement, jusqu'à obtenir $|V|$ unités. Un mot courant est un seul token ; un mot rare est
découpé en morceaux.

\begin{exemple}[(fusions sur un mini-corpus)]
Corpus : « vol » (5 fois), « volé » (3), « vélo » (2), découpés en lettres.
Paire la plus fréquente : (v,o) dans « vol » et « volé » : $5+3=8$ $\Rightarrow$ token « vo ».
Puis (vo,l) : 8 $\Rightarrow$ « vol ». Puis (vol,é) : 3 $\Rightarrow$ « volé ». « vélo » reste
découpé en v, é, l, o. Chaque token reçoit un numéro, son \emph{identifiant}.
\end{exemple}

\subsection{De l'identifiant au vecteur}
Une séquence de $N$ tokens d'identifiants $v_1,\dots,v_N$ devient la matrice
$\mat{H}^{(0)}\in\R^{N\times d}$ dont la ligne $s$ est la ligne $v_s$ de $\mat{E}$.

\subsection{Un bloc transformer}
Chacune des $L=28$ couches transforme $\mat{H}^{(\ell-1)}$ en $\mat{H}^{(\ell)}$ :
\[
\mat{A} = \mat{H}^{(\ell-1)} + \mathrm{Attn}\big(\mathrm{RMSNorm}(\mat{H}^{(\ell-1)})\big),\qquad
\mat{H}^{(\ell)} = \mat{A} + \mathrm{FFN}\big(\mathrm{RMSNorm}(\mat{A})\big).
\]
\textbf{RMSNorm} d'une ligne $\vect{h}\in\R^d$ :
$\mathrm{RMSNorm}(\vect{h}) = \dfrac{\vect{h}}{\sqrt{\frac1d\sum_{u=1}^{d}h_u^2+\varepsilon}}\odot\vect{g}$,
avec $\varepsilon=10^{-5}$ et $\vect{g}$ un vecteur de gains appris ($\odot$ : produit terme à terme).

\textbf{Attention causale} (une tête) :
\[
\mat{Q}=\mat{H}\mat{W}_Q,\quad \mat{K}=\mat{H}\mat{W}_K,\quad \mat{V}=\mat{H}\mat{W}_V,\qquad
\mathrm{Attn}(\mat{H}) = \softmax\!\Big(\frac{\mat{Q}\mat{K}^\top}{\sqrt{d_h}} + \mat{M}\Big)\mat{V},
\]
où le softmax est appliqué ligne par ligne et $\mat{M}$ est le masque causal : $M_{s,s'}=0$ si
$s'\le s$ et $-\infty$ sinon (un token ne voit pas les suivants). Le coefficient $(s,s')$ de la
matrice softmax est le poids que le token $s$ accorde au token $s'$.

\begin{exemple}[(attention sur 3 tokens, $d_h = 2$)]
$\mat{H}=\begin{pmatrix}1&0\\0&1\\1&1\end{pmatrix}$,
$\mat{W}_Q=\begin{pmatrix}1&0\\0&1\end{pmatrix}$,
$\mat{W}_K=\begin{pmatrix}1&1\\0&1\end{pmatrix}$,
$\mat{W}_V=\begin{pmatrix}1&0\\1&1\end{pmatrix}$.

$\mat{Q}=\begin{pmatrix}1&0\\0&1\\1&1\end{pmatrix}$,
$\mat{K}=\begin{pmatrix}1&1\\0&1\\1&2\end{pmatrix}$,
$\mat{V}=\begin{pmatrix}1&0\\1&1\\2&1\end{pmatrix}$,
$\dfrac{\mat{Q}\mat{K}^\top}{\sqrt2}=\begin{pmatrix}0{,}707&0&0{,}707\\0{,}707&0{,}707&1{,}414\\1{,}414&0{,}707&2{,}121\end{pmatrix}$.

Après masque et softmax par ligne :
$\begin{pmatrix}1&0&0\\0{,}5&0{,}5&0\\0{,}284&0{,}140&0{,}576\end{pmatrix}$ ; ligne 3 :
$e^{1{,}414},e^{0{,}707},e^{2{,}121}$ normalisés.
Sortie : $\begin{pmatrix}1&0\\1&0{,}5\\1{,}576&0{,}716\end{pmatrix}$ : chaque token devient une
moyenne pondérée des valeurs des tokens qu'il peut voir.
\end{exemple}

\textbf{Plusieurs têtes et GQA.} Les $H=24$ têtes calculent ces opérations en parallèle sur des
projections de dimension $d_h=128$ ; leurs sorties sont concaténées ($24\times128=3\,072$) puis
multipliées par $\mat{W}_O$. Avec l'attention à requêtes groupées (GQA), 3 têtes de requêtes
partagent la même paire $(\mat{K},\mat{V})$ : seules 8 paires sont calculées et stockées.

\textbf{Position (RoPE).} La position $s$ d'un token est introduite en faisant tourner chaque paire de
coordonnées $(q_{2u},q_{2u+1})$ de $\vect{q}_s$ et $\vect{k}_s$ d'un angle $s\,\theta_u$, avec
$\theta_u = 500\,000^{-2u/d_h}$ ; le produit scalaire $\vect{q}_s\cdot\vect{k}_{s'}$ ne dépend alors
que de l'écart $s-s'$.

\textbf{Réseau feed-forward SwiGLU :}
$\mathrm{FFN}(\vect{h}) = \mat{W}_2\big(\mathrm{SiLU}(\mat{W}_1\vect{h})\odot \mat{W}_3\vect{h}\big)$, avec
$\mat{W}_1,\mat{W}_3\in\R^{8\,192\times3\,072}$, $\mat{W}_2\in\R^{3\,072\times8\,192}$ et
$\mathrm{SiLU}(a)=a/(1+e^{-a})$.

\subsection{Génération et réponse JSON imposée}
La dernière ligne $\vect{h}_N$ (après une RMSNorm) donne les logits du token suivant :
$\vect{z} = \mat{E}\,\vect{h}_N \in \R^{|V|}$. Avec une température $\tau$,
$P(v) = \softmax(\vect{z}/\tau)_v$. L'application fixe $\tau=0$ : on prend
$v^\star = \argmax_v z_v$ (décodage glouton), donc la même question donne la même réponse.

Le paramètre \texttt{format} d'Ollama reçoit un schéma JSON. Il est traduit en grammaire, et à chaque
pas les tokens qui rendraient la sortie invalide reçoivent un logit $-\infty$ :
$z'_v = z_v + m_v$, avec $m_v = 0$ si $v$ est autorisé et $-\infty$ sinon.

\begin{exemple}[(masque de grammaire)]
Après \texttt{\{"reponse": "}, seuls « oui », « non », « incertain » sont permis. Logits bruts :
« peut-être » 3,1 ; « oui » 2,4 ; « non » 1,9. Sans masque, le modèle écrirait « peut-être » ;
avec le masque, il écrit « oui ». La forme est garantie ; le fond reste à vérifier.
\end{exemple}

\begin{algorithm}[h]
\caption{Boucle de génération (Ollama, $\tau=0$)}
\KwIn{prompt tokenisé $v_1..v_N$, schéma JSON, \texttt{num\_predict}=400}
calculer les couches sur $v_1..v_N$ et garder $\mat{K},\mat{V}$ de chaque couche (cache KV)\;
\Repeat{token de fin, JSON complet ou 400 tokens}{
  $\vect{z}\leftarrow$ logits du dernier token ; $\vect{z}\leftarrow\vect{z}+\vect{m}$ (masque de grammaire)\;
  $v\leftarrow\argmax\vect{z}$ ; ajouter $v$ ; calculer sa ligne dans chaque couche en réutilisant le cache\;
}
\Return texte décodé\;
\end{algorithm}

\begin{exemple}[(mémoire du cache KV, \texttt{num\_ctx} = 4\,096)]
$2\ (\mat{K}\text{ et }\mat{V})\times 28\ \text{couches}\times 8\ \text{têtes}\times128\times4\,096\ \text{positions}\times2\ \text{octets} = 469\,762\,048$ octets
$= 448$ Mio. Avec les poids (1,95 Go), le modèle tient dans environ 2,5 Go, ce qui explique
qu'il tourne sans GPU, mais lentement : chaque token demande environ $2\times3{,}2\times10^9$
multiplications-additions.
\end{exemple}

\subsection{Les trois questions fermées et leur vérification}
Pour les contrôles 4 (événement et bien garantis), 5 (exclusions) et 6 (conditions), le modèle
reçoit les seules clauses du type utile et doit répondre « oui », « non » ou « incertain », en
citant une clause $C_r$ et un extrait. L'application accepte la réponse seulement si l'extrait
$e$ figure dans la clause $c$ : inclusion exacte après normalisation, ou recouvrement
\[
\rho(e,c) = \frac{1}{|e|}\sum_{w \in e}\mathbb{1}[w \in \mathrm{mots}(c)] \ \ge\ 0{,}8 ,
\qquad |e|\ge 3 ,
\]
où $|e|$ est le nombre de mots de l'extrait. Sinon le contrôle reste « à vérifier ».

\begin{exemple}[(exclusion Sogessur)]
Clause : « \dots les vols commis dans les locaux à usage commun des occupants \dots »
\begin{itemize}
\item « les vols commis dans les locaux à usage commun des occupants » : inclusion exacte $\Rightarrow$ acceptée ;
\item « vols dans les locaux communs des occupants de l immeuble » : 6 mots sur 10 dans la clause,
$\rho=0{,}6<0{,}8$ $\Rightarrow$ refusée (« communs » et « immeuble » ne sont pas dans la clause).
\end{itemize}
\end{exemple}

%=====================================================================
\section{L'instruction en 7 contrôles et le verdict}

Chaque contrôle $c=1,\dots,7$ reçoit un statut $\sigma_c\in\{\text{satisfait},\text{bloquant},
\text{à vérifier},\text{sans objet}\}$. Les contrôles 1, 2, 3, 6 (délai) et 7 sont factuels
(requêtes SQL et calculs) ; 4, 5 et 6 (conditions) sont interprétés par le modèle.

\textbf{Agrégation non compensatoire :}
\[
\text{verdict} =
\begin{cases}
\text{potentiellement exclu} & \text{si } \exists\, c:\ \sigma_c = \text{bloquant},\\
\text{potentiellement couvert} & \text{si } \forall\, c:\ \sigma_c \in\{\text{satisfait},\text{sans objet}\},\\
\text{indéterminé} & \text{sinon.}
\end{cases}
\]
Un contrôle favorable ne compense jamais un contrôle bloquant : c'est la logique d'un gestionnaire
(une exclusion écarte la garantie quel que soit le reste).

\textbf{Garantie acquise} (contrôle 3) : une garantie $g$ est acquise par la police $\pi$ si
$g$ est de base, ou si la formule de $\pi$ figure parmi les formules de $g$, ou si l'option de
$g$ fait partie des options de $\pi$ ; si le bien touché relève d'une extension (pack jardin), cette
extension doit aussi être acquise.

\textbf{Délai de déclaration} (contrôle 6) : en jours calendaires $\Delta = t_{\text{décl}}-t_{\text{conn}}$,
ou en jours ouvrés $\Delta = \sum_{u=t_{\text{conn}}+1}^{t_{\text{décl}}}\mathbb{1}[u \text{ est un jour de lundi à vendredi}]$.
Un retard n'est jamais bloquant seul : le droit luxembourgeois, belge et français exige un
préjudice de l'assureur.

\textbf{Montant} (contrôle 7) : $I = \min\big(\max(D - F,\,0),\ P\big)$, avec $D$ le dommage
estimé, $F$ la franchise et $P$ le plus petit plafond applicable (avant vétusté).

\begin{exemple}[(sinistres de démonstration)]
S-02 (vélo, Ethias, option vol) : $D=900$, $F=300$, $P=45\,000$ :
$I=\min(\max(600,0),45\,000)=600$ €. Déclaré 3 jours après la connaissance, délai 10 jours :
satisfait.

S-11 (bougie, Baloise, garanties de base) : la garantie incendie est acquise, mais le mobilier de
jardin n'est couvert que par le « Pack jardin / piscine », non souscrit : contrôle 3 bloquant,
donc verdict « potentiellement exclu (garantie non souscrite) », sans appel au modèle.
\end{exemple}

%=====================================================================
\section{Mesurer : protocole et statistiques}

\subsection{Protocole}
Le jeu \texttt{cas\_couverture.csv} contient 13 sinistres $\times$ 3 contrats $=39$ jugements
$(i,j)$, répartis à l'avance en dev (18) et test (21). Tout réglage se fait sur dev ; la mesure sur
test est faite une seule fois. Chaque jugement indique le verdict attendu, les verdicts acceptables et
un fragment de la clause décisive retrouvé mot pour mot. Mesures :
verdict exact $\frac{1}{N}\sum \mathbb{1}[\hat v = v]$, fausse couverture (verdict « couvert » sur
un cas non couvert), faux refus, et rappel de la clause décisive transmise au modèle.

\textbf{Référence.} Sans modèle, les contrôles 4 à 6 restent « à vérifier » : 29\,\% de verdicts
exacts sur test, aucune fausse couverture. Toute amélioration se mesure contre cette référence.

\subsection{Intervalle de Wilson}
Pour $k$ succès sur $N$ essais, $\hat p = k/N$, et $z=1{,}96$ pour 95\,\% :
\[
\frac{\hat p + \frac{z^2}{2N}}{1+\frac{z^2}{N}} \ \pm\
\frac{z}{1+\frac{z^2}{N}}\sqrt{\frac{\hat p(1-\hat p)}{N}+\frac{z^2}{4N^2}} .
\]
Contrairement à l'intervalle « $\hat p\pm z\sqrt{\hat p(1-\hat p)/N}$ », il reste dans $[0,1]$ et
reste fiable pour de petits $N$ ou des proportions proches de 0 ou 1.

\begin{exemple}[($k=24$, $N=28$)]
$\hat p = 0{,}857$ ; centre $=(0{,}857+0{,}0686)/(1+0{,}1372)=0{,}814$ ;
demi-largeur $=1{,}96/1{,}1372\times\sqrt{0{,}00437+0{,}00122}=0{,}129$ ; intervalle
$[68{,}5\,\% ;\ 94{,}3\,\%]$.
\end{exemple}

\subsection{Test de McNemar exact}
Pour comparer deux versions A et B sur les mêmes cas, seuls comptent les cas discordants :
$b$ (B réussit, A échoue) et $c$ (l'inverse). Sous l'hypothèse « aucune différence », chaque
discordance tombe d'un côté ou de l'autre avec probabilité $\frac12$ :
\[
p = \min\Big(1,\ 2\sum_{u=0}^{\min(b,c)}\binom{b+c}{u}2^{-(b+c)}\Big).
\]
\begin{exemple}[(projet 1, 19.4.0 contre 19.5.0)]
$b=8$, $c=2$ : $p = 2\,\big(\binom{10}{0}+\binom{10}{1}+\binom{10}{2}\big)/2^{10} = 2\times56/1\,024 = 0{,}109$.
Le gain n'est pas significatif au seuil de 5\,\%.
\end{exemple}

\subsection{Bootstrap}
Pour un indicateur sans formule d'intervalle (F1 macro), on tire 2\,000 fois, avec remise, $m$ cas
parmi les $m$ cas de test, on recalcule l'indicateur, et l'on garde les quantiles 2,5\,\% et 97,5\,\%.

%=====================================================================
\section{Rappels : Insurance Claims Assistance}

\textbf{Similarité cosinus (bge-m3).} Question $q_i$ et document $d_j$ sont encodés en
$\vect{u}_i,\vect{v}_j\in\R^{1\,024}$ ; $\cos(\vect{u}_i,\vect{v}_j) = \dfrac{\vect{u}_i\cdot\vect{v}_j}{\|\vect{u}_i\|\,\|\vect{v}_j\|}$.
Pour les 201 documents, c'est un produit matrice-vecteur : $\mat{D}\vect{u}_i$, avec $\mat{D}\in\R^{201\times1\,024}$
dont les lignes sont normalisées. Exemple : $(0{,}2;0{,}1;0{,}9)$ et $(0{,}3;0;0{,}8)$ donnent $0{,}984$.

\textbf{Fusion des rangs (RRF).} $\mathrm{RRF}(d_j)=\sum_{s}\frac{1}{k+\mathrm{rang}_s(d_j)}$, $k=60$.
Un document 1\textsuperscript{er} et 3\textsuperscript{e} : $1/61+1/63=0{,}03227$ ; un document 2\textsuperscript{e}
et 1\textsuperscript{er} : $0{,}03252$, classé devant.

\textbf{Justification par inférence textuelle (essai non retenu).}
$\mathrm{score} = \min_{h}\max_{p} P(\text{implication}\mid p,h)$, où $h$ parcourt les phrases de la
réponse et $p$ les prémisses (extraits ou fenêtres de 3 phrases). Bloquer sous un seuil de 0,95
faisait passer les refus corrects de 5/28 à 24/28, mais les réponses conformes de 15/30 à 8/30
($c=7$, $b=0$, $p=0{,}016$) : contrôle non intégré.

%=====================================================================
\section{Limites connues}
\begin{itemize}
\item Le classifieur n'a pas de classe « hors habitation » : une question sur l'assurance décès est
rangée dans une famille habitation (mesuré sur dev).
\item La clause « dommages intentionnels » n'est pas transmise au modèle pour l'incendie volontaire
(mesuré sur test, donc non corrigé dans cette version).
\item Les jours fériés ne sont pas déduits des jours ouvrés.
\item Les verdicts attendus sont établis par lecture des conditions générales ; ils doivent être
revus par un second lecteur métier.
\item Les intervalles supposent des passages indépendants ; les passages répétés d'un même cas ne
le sont pas, ce qui rend les intervalles légèrement optimistes.
\end{itemize}

%=====================================================================
\section*{Références}
\addcontentsline{toc}{section}{Références}
\small
\begin{enumerate}[label={[\arabic*]}]
\item Manning, C. D., Raghavan, P., Schütze, H. (2008). \emph{Introduction to Information Retrieval}. Cambridge University Press, chap. 6 (TF-IDF) et 11 (modèles probabilistes). En ligne : nlp.stanford.edu/IR-book.
\item Robertson, S., Zaragoza, H. (2009). The Probabilistic Relevance Framework: BM25 and Beyond. \emph{Foundations and Trends in Information Retrieval}, 3(4), 333--389.
\item Spärck Jones, K. (1972). A statistical interpretation of term specificity and its application in retrieval. \emph{Journal of Documentation}, 28(1), 11--21.
\item SQLite. \emph{FTS5 Extension}, section « The bm25() function ». sqlite.org/fts5.html.
\item Bishop, C. M. (2006). \emph{Pattern Recognition and Machine Learning}. Springer, chap. 4.3 (régression logistique multinomiale).
\item Hastie, T., Tibshirani, R., Friedman, J. (2009). \emph{The Elements of Statistical Learning}, 2\textsuperscript{e} éd. Springer, chap. 4 et 7 (bootstrap).
\item Liu, D. C., Nocedal, J. (1989). On the limited memory BFGS method for large scale optimization. \emph{Mathematical Programming}, 45, 503--528.
\item Pedregosa, F. et al. (2011). Scikit-learn: Machine Learning in Python. \emph{JMLR}, 12, 2825--2830 ; documentation de \texttt{TfidfVectorizer}.
\item Vaswani, A. et al. (2017). Attention Is All You Need. \emph{NeurIPS}. arXiv:1706.03762.
\item Grattafiori, A. et al. (2024). The Llama 3 Herd of Models. arXiv:2407.21783.
\item Meta (2024). Llama 3.2 : fiche du modèle et fichier \texttt{config.json} du modèle 3B (Hugging Face, meta-llama/Llama-3.2-3B-Instruct).
\item Sennrich, R., Haddow, B., Birch, A. (2016). Neural Machine Translation of Rare Words with Subword Units. \emph{ACL}. arXiv:1508.07909.
\item Zhang, B., Sennrich, R. (2019). Root Mean Square Layer Normalization. \emph{NeurIPS}. arXiv:1910.07467.
\item Su, J. et al. (2021). RoFormer: Enhanced Transformer with Rotary Position Embedding. arXiv:2104.09864.
\item Shazeer, N. (2020). GLU Variants Improve Transformer. arXiv:2002.05202.
\item Ainslie, J. et al. (2023). GQA: Training Generalized Multi-Query Transformer Models from Multi-Head Checkpoints. \emph{EMNLP}. arXiv:2305.13245.
\item Willard, B. T., Louf, R. (2023). Efficient Guided Generation for Large Language Models. arXiv:2307.09702.
\item Gerganov, G. et al. \emph{llama.cpp} et format GGUF (quantification Q4\_K\_M). github.com/ggml-org/llama.cpp ; Ollama, documentation de l'API (\texttt{/api/generate}, paramètre \texttt{format}). github.com/ollama/ollama.
\item Jurafsky, D., Martin, J. H. (2025). \emph{Speech and Language Processing}, 3\textsuperscript{e} éd. (brouillon en ligne), chap. 9 à 11 (transformers, LLM).
\item Wilson, E. B. (1927). Probable Inference, the Law of Succession, and Statistical Inference. \emph{JASA}, 22(158), 209--212.
\item Brown, L. D., Cai, T. T., DasGupta, A. (2001). Interval Estimation for a Binomial Proportion. \emph{Statistical Science}, 16(2), 101--133.
\item McNemar, Q. (1947). Note on the sampling error of the difference between correlated proportions or percentages. \emph{Psychometrika}, 12(2), 153--157.
\item Agresti, A. (2013). \emph{Categorical Data Analysis}, 3\textsuperscript{e} éd. Wiley (tests appariés, intervalles).
\item Efron, B., Tibshirani, R. (1993). \emph{An Introduction to the Bootstrap}. Chapman \& Hall.
\item Brier, G. W. (1950). Verification of forecasts expressed in terms of probability. \emph{Monthly Weather Review}, 78(1), 1--3.
\item Chen, J. et al. (2024). BGE M3-Embedding. arXiv:2402.03216 ; Cormack, G. V., Clarke, C., Büttcher, S. (2009). Reciprocal Rank Fusion. \emph{SIGIR}.
\item He, P. et al. (2021). DeBERTa: Decoding-enhanced BERT with Disentangled Attention. \emph{ICLR} ; Williams, A., Nangia, N., Bowman, S. (2018). A Broad-Coverage Challenge Corpus for Sentence Understanding through Inference (MNLI). \emph{NAACL}.
\end{enumerate}
\end{document}
